% !TEX program = xelatex
\documentclass[10pt,a4paper]{article}

% ---- Пакеты для XeLaTeX и кириллицы ----
\usepackage{fontspec}
\usepackage{polyglossia}
\setmainlanguage{russian}
\setotherlanguage{english}

% ---- Шрифт с кириллицей ----
\setmainfont{Liberation Serif}
\setsansfont{Liberation Sans}
\setmonofont{Liberation Mono}

% ---- Поля и геометрия ----
\usepackage[top=1.5cm, bottom=1.5cm, left=1.5cm, right=1.5cm]{geometry}
\usepackage{parskip}
\setlength{\parindent}{0pt}

% ---- Продвинутые ссылки ----
\usepackage{hyperref}
\hypersetup{
    colorlinks=true,
    linkcolor=blue,
    urlcolor=blue,
    citecolor=blue
}

% ---- Полезные пакеты ----
\usepackage{enumitem}
\setlist{nosep, left=0pt}
\usepackage{tabularx}
\usepackage{array}
\usepackage{booktabs}
\usepackage{xcolor}
\usepackage{titlesec}

% ---- Настройка заголовков разделов ----
\titleformat{\section}{\Large\bfseries\uppercase}{}{0em}{}[\vspace{-0.3em}\rule{\textwidth}{0.5pt}]
\titlespacing*{\section}{0pt}{0.8em}{0.3em}

\titleformat{\subsection}{\large\bfseries}{}{0em}{}
\titlespacing*{\subsection}{0pt}{0.4em}{0.2em}

% ---- Пользовательские команды ----
\newcommand{\projectEntry}[3]{\textbf{#1} \hfill \textit{#2}\par #3\par\vspace{0.3em}}

% ---- Общие данные (имя, телефон, email и т.д.) ----
% ============================================================
% ОБЩИЕ ПЕРСОНАЛЬНЫЕ ДАННЫЕ
% ============================================================
% Меняешь здесь — меняется во всех резюме

\newcommand{\PersonName}{Кузык Иван Романович}
\newcommand{\PersonPhone}{+7 (916) 574-52-99}
\newcommand{\PersonEmail}{vanya.kuzyk1@gmail.com}
\newcommand{\PersonGitHub}{github.com/vanya890}
\newcommand{\PersonCity}{г.~Москва, м.~Бульвар Дмитрия Донского}
\newcommand{\PersonCitizenship}{РФ}
\newcommand{\PersonDrivingLicense}{Категория B}

% ============================================================
% КОМАНДА ДЛЯ ЗАГОЛОВКА РЕЗЮМЕ
% ============================================================
% Использование: \makeHeader{Должность}
\newcommand{\makeHeader}[1]{%
\begin{center}
    {\Huge\bfseries \PersonName}\\[0.3em]
    {\large #1}\\[0.5em]
    \begin{tabular}{ll}
        Телефон: & \PersonPhone \\
        Email:   & \href{mailto:\PersonEmail}{\PersonEmail} \\
        GitHub:  & \href{https://\PersonGitHub}{\PersonGitHub} \\
        Город:   & \PersonCity \\
        Гражданство: & \PersonCitizenship
    \end{tabular}
\end{center}
}

% ============================================================
% ОБЩЕЕ ОБРАЗОВАНИЕ
% ============================================================
% Использование: \makeEducation
\newcommand{\makeEducation}{%
\section{Образование}

\projectEntry{Московский политехнический университет --- Факультет информационных технологий}{2022--2026}
{Бакалавриат, направление «Информатика и вычислительная техника», специализация «Киберфизические системы». ВКР «Разработка образовательного 3D-симулятора робототехники» --- отлично. Средний балл --- 5.0.}

\projectEntry{ГБПОУ «Образовательный комплекс «Юго-Запад»}{2018--2022}
{Среднее специальное образование, специальность «Информационные системы и программирование». Диплом с отличием. Демонстрационный экзамен по компетенции «Разработка компьютерных игр и мультимедийных приложений» --- отлично.}
}

% ============================================================
% ОБЩАЯ ДОПОЛНИТЕЛЬНАЯ ИНФОРМАЦИЯ
% ============================================================
% Использование: \makeFooter{Ключевые компетенции}{Личные качества}
\newcommand{\makeFooter}[2]{%
\section{Дополнительная информация}

\begin{itemize}
    \item \textbf{GitHub:} \href{https://\PersonGitHub}{\PersonGitHub} --- #1
    \item \textbf{Ключевые компетенции:} #2
    \item \textbf{Личные качества:} Обучаемость, ответственность, умение работать с большим объёмом информации, системное мышление, ориентация на результат
    \item \textbf{Водительские права:} \PersonDrivingLicense
\end{itemize}
}

\begin{document}

% ============================================================
% ЗАГОЛОВОК
% ============================================================
\makeHeader{Аналитик данных / Data Scientist (Junior)}

% ============================================================
% О СЕБЕ
% ============================================================
\section{О себе}

Аналитик данных с техническим образованием в области информатики и вычислительной техники (Московский политехнический университет). Владею \textbf{Python (NumPy, Pandas, Matplotlib, SciPy)}, методами статистического анализа, машинного обучения и визуализации данных. Имею практический опыт разработки нейросетевых моделей, численных методов, корреляционного, факторного и кластерного анализа. Умею работать с базами данных (SQL), обрабатывать и очищать данные, строить дашборды и презентовать результаты. Ориентирован на решение бизнес-задач с помощью данных.

% ============================================================
% ОПЫТ РАБОТЫ (ПРОЕКТНАЯ ДЕЯТЕЛЬНОСТЬ)
% ============================================================
\section{Опыт работы}

\subsection{Проектная деятельность / Фриланс}

\projectEntry{Аналитик данных (Курсовая: Программная обработка данных мониторинга)}{2025–2026}
{Разработал desktop-приложение на \textbf{C\# Windows Forms} для статистической обработки данных мониторинга. Реализовал вручную (без сторонних библиотек):
\begin{itemize}
    \item Расчёт описательной статистики, исключение выбросов, определение квалификационных уровней
    \item Корреляционный анализ (матрица корреляций Пирсона)
    \item Факторный анализ (собственные значения и векторы)
    \item Кластерный анализ с построением дендрограммы
    \item Анализ качества теста: распределение заданий по сложности и дифференцирующей способности
    \item Оптимизация теста на основе анализа заданий
\end{itemize}}

\projectEntry{Исследователь (AI\_Study\_Colab — нейросети и ML)}{2024–2026}
{Выполнил цикл лабораторных и контрольных работ по искусственному интеллекту и машинному обучению в \textbf{Jupyter Notebook / Google Colab}:
\begin{itemize}
    \item Реализовал модель нейрона и алгоритмы обучения
    \item Создал и обучил нейросеть для задач классификации
    \item Работал с библиотеками \textbf{NumPy, Matplotlib} для численных расчётов и визуализации
    \item Выполнил задания по синтаксису Python, анализу данных, построению графиков
\end{itemize}}

\projectEntry{Численные методы (NumberMethods)}{2024}
{Реализовал полный набор классических алгоритмов численного анализа на \textbf{Python}:
\begin{itemize}
    \item Решение СЛАУ методом Гаусса с выбором главного элемента
    \item Интерполяция и аппроксимация функций
    \item Численное дифференцирование и интегрирование
    \item Визуализация результатов с помощью Matplotlib
\end{itemize}}

\projectEntry{Аналитик данных (Stroy-Materiali-24.ru — SEO-аналитика)}{2024}
{Разработал инструменты для SEO-оптимизации сайта строительных материалов. Реализовал:
\begin{itemize}
    \item Парсинг и анализ популярных поисковых запросов (CSV)
    \item Обработку Markdown-контента с помощью Python-скриптов
    \item Генерацию контента с использованием AI-промптов
    \item Анализ данных вебмастера для оптимизации ключевых слов
\end{itemize}}

\projectEntry{AI Chat Game — аналитика токенов}{2024–2025}
{Разработал дашборд аналитики токенов для веб-приложения AI-чата:
\begin{itemize}
    \item Ежедневная/еженедельная/ежемесячная статистика использования токенов
    \item Разбивка затрат по моделям (gpt-4o-mini vs gpt-4o)
    \item Рейтинг персонажей по использованию
    \item Прогнозы расходов и рекомендации по оптимизации
    \item Визуализация с помощью Recharts
\end{itemize}}

\subsection{Практика}

\projectEntry{Оператор БД (Издательство «Просвещение»)}{Май–Июль 2021}
{Заполнение базы данных информацией, описание изображений для поиска по ключевым словам, подготовка данных для редакций.}

\projectEntry{Оператор ПК (ГБПОУ «Юго-Запад»)}{Ноябрь 2020 – Январь 2021}
{Ввод и редактирование данных в программе \textbf{1С:Библиотека}, обучение сотрудников работе с системой на базовом уровне.}

% ============================================================
% ОБРАЗОВАНИЕ
% ============================================================
\makeEducation

% ============================================================
% КЛЮЧЕВЫЕ ДИСЦИПЛИНЫ
% ============================================================
\section{Ключевые дисциплины}

\begin{tabular}{p{6cm} p{10cm}}
    \textbf{Математика и статистика:} & Математический анализ, теория вероятностей и мат. статистика, дискретная математика, численные методы \\
    \textbf{Анализ данных:} & Теория систем и системный анализ, программная обработка данных, методы и алгоритмы управления \\
    \textbf{AI/ML:} & Искусственные нейронные сети, нейронные сети глубокого обучения в обработке изображений, интеллектуальные системы \\
    \textbf{Базы данных:} & Основы проектирования баз данных, администрирование и автоматизация БД, SQL
\end{tabular}

% ============================================================
% НАВЫКИ
% ============================================================
\section{Навыки}

\begin{tabular}{p{3.5cm} p{12cm}}
    \textbf{Языки:} & Python, SQL, C\#, TypeScript, JavaScript \\
    \textbf{Анализ данных:} & NumPy, Pandas, SciPy, Matplotlib, статистический анализ, корреляционный анализ, факторный анализ, кластерный анализ \\
    \textbf{Machine Learning:} & Нейронные сети (CNN), модели нейронов, обучение нейросетей, Scikit-learn (базовый) \\
    \textbf{Базы данных:} & MySQL, SQLite, Oracle (базовый), 1С \\
    \textbf{Инструменты:} & Git, Jupyter Notebook, Google Colab, Excel \\
    \textbf{Визуализация:} & Matplotlib, Recharts, построение дашбордов \\
    \textbf{Прочее:} & Docker, Linux, REST API, JSON, CSV \\
    \textbf{Языки:} & Русский — родной, Английский — A2 (элементарный)
\end{tabular}

% ============================================================
% ДОПОЛНИТЕЛЬНАЯ ИНФОРМАЦИЯ
% ============================================================
\makeFooter{репозитории с реализациями алгоритмов, нейросетей, численных методов}{Сбор и очистка данных, статистический анализ, визуализация, машинное обучение, работа с базами данных, написание отчётов}

\end{document}